\section{Introduction}
\label{sec:research_objectives_introduction}

The background chapter reviewed current approaches to compositional generative
modelling and highlighted key challenges that remain unresolved in the
literature. Score-based methods (e.g., Liu et al.) linearly combine diffusion
scores but often fail when the underlying generative factors are entangled or
misaligned, producing compositions that do not correspond to meaningful semantic
transformations. Energy-based formulations (e.g., Du et al.) cast composition as
a Product-of-Experts, but rely on strong independence or factorisation
assumptions; when these assumptions do not hold, PoE models commonly collapse or
amplify incompatible modes.

More recently, density-based approaches such as SuperDiff provide a
mathematically principled way to estimate and combine log densities along a
shared diffusion trajectory. However, while these methods guarantee correct
density computation under shared forward dynamics, they do \emph{not} guarantee
that density-level operations correspond to meaningful semantic combinations.
Their limitations are therefore conceptual rather than algorithmic: they offer a
precise way to combine distributions, but lack an underlying theory explaining
when such combinations reflect the true semantic factors of the domain.

Motivated by these gaps, our research seeks to (i) clarify when and why
pretrained generative models yield semantically meaningful composites,
(ii) extend compositional methods to robust property-guided settings, and
(iii) identify the structural and geometric preconditions for compositional
correctness. These goals form the basis of the research objectives and
hypotheses formalised in the remainder of this chapter.
